\documentclass[10pt,a4paper]{amsart}
\usepackage{amsmath,amssymb,amsthm,graphicx}
\usepackage{verbatim}
\newenvironment{reference}{\comment}{\endcomment}
%\newenvironment{reference}{}{}
\newenvironment{slogan}{\comment}{\endcomment}
\newenvironment{history}{\comment}{\endcomment}

% For commutative diagrams we use Xy-pic
\usepackage[all]{xy}

% We use 2cell for 2-commutative diagrams.
\xyoption{2cell}
% Theorem environments.
%
\theoremstyle{plain}
\newtheorem{theorem}[subsection]{Theorem}
\newtheorem{proposition}[subsection]{Proposition}
\newtheorem{lemma}[subsection]{Lemma}

\theoremstyle{definition}
\newtheorem{definition}[subsection]{Definition}
\newtheorem{example}[subsection]{Example}
\newtheorem{exercise}[subsection]{Exercise}
\newtheorem{situation}[subsection]{Situation}

\theoremstyle{remark}
\newtheorem{remark}[subsection]{Remark}
\newtheorem{remarks}[subsection]{Remarks}

\numberwithin{equation}{subsection}

% Macros
%
\def\lim{\mathop{\mathrm{lim}}\nolimits}
\def\colim{\mathop{\mathrm{colim}}\nolimits}
\def\Spec{\mathop{\mathrm{Spec}}}
\def\Hom{\mathop{\mathrm{Hom}}\nolimits}
\def\Ext{\mathop{\mathrm{Ext}}\nolimits}
\def\SheafHom{\mathop{\mathcal{H}\!\mathit{om}}\nolimits}
\def\SheafExt{\mathop{\mathcal{E}\!\mathit{xt}}\nolimits}
\def\Sch{\mathit{Sch}}
\def\Mor{\mathop{\mathrm{Mor}}\nolimits}
\def\Ob{\mathop{\mathrm{Ob}}\nolimits}
\def\Sh{\mathop{\mathit{Sh}}\nolimits}
\def\NL{\mathop{N\!L}\nolimits}
\def\CH{\mathop{\mathrm{CH}}\nolimits}
\def\proetale{{pro\text{-}\acute{e}tale}}
\def\etale{{\acute{e}tale}}
\def\QCoh{\mathit{QCoh}}
\def\Ker{\mathop{\mathrm{Ker}}}
\def\Im{\mathop{\mathrm{Im}}}
\def\Coker{\mathop{\mathrm{Coker}}}
\def\Coim{\mathop{\mathrm{Coim}}}

% Boxtimes
%
\DeclareMathSymbol{\boxtimes}{\mathbin}{AMSa}{"02}

%
% Macros for moduli stacks/spaces
%
\def\QCohstack{\mathcal{QC}\!\mathit{oh}}
\def\Cohstack{\mathcal{C}\!\mathit{oh}}
\def\Spacesstack{\mathcal{S}\!\mathit{paces}}
\def\Quotfunctor{\mathrm{Quot}}
\def\Hilbfunctor{\mathrm{Hilb}}
\def\Curvesstack{\mathcal{C}\!\mathit{urves}}
\def\Polarizedstack{\mathcal{P}\!\mathit{olarized}}
\def\Complexesstack{\mathcal{C}\!\mathit{omplexes}}
% \Pic is the operator that assigns to X its picard group, usage \Pic(X)
% \Picardstack_{X/B} denotes the Picard stack of X over B
% \Picardfunctor_{X/B} denotes the Picard functor of X over B
\def\Pic{\mathop{\mathrm{Pic}}\nolimits}
\def\Picardstack{\mathcal{P}\!\mathit{ic}}
\def\Picardfunctor{\mathrm{Pic}}
\def\Deformationcategory{\mathcal{D}\!\mathit{ef}}
\usepackage{cleveref}
\crefname{originallemma}{original lemma}{original lemmas}
\Crefname{originallemma}{Original lemma}{Original lemmas}
\makeatletter
\newlabel{lemma-Zariski-topology}{{00E0}{source}{}{sourcepointer-1}{}}
\newlabel{lemma-Zariski-topology@cref}{{[originallemma][0][]00E0}{[1][0][]source}{}{}{}}
\makeatother
\crefname{originallemma}{original lemma}{original lemmas}
\Crefname{originallemma}{Original lemma}{Original lemmas}
\makeatletter
\newlabel{lemma-in-image}{{00E7}{source}{}{sourcepointer-2}{}}
\newlabel{lemma-in-image@cref}{{[originallemma][0][]00E7}{[1][0][]source}{}{}{}}
\makeatother
\title{Surjectivity of the map on spectra is equivalent to the contraction of extended ideals being the ideals themselves}
\author{The Stacks Project Authors}
\date{stacks-project:algebra.tex:local-snapshot-20261012; GNU FDL 1.2 or later}
\begin{document}
\maketitle
\noindent\textit{Editorial note.} All delivered mathematical text is the author's own, copied from the pinned Stacks Project chapter file algebra.tex; no statement or proof step was rewritten, reordered or removed.\par\medskip
\noindent\textit{Editorial note.} Preamble provenance, as required for this source: the frozen snapshot stages the chapter's own preamble file beside it (preamble.tex, 217 lines, 8116 bytes, the same shared Stacks-wide preamble the published chow and etale-cohomology cards use); the first line of the pinned chapter file is \textbackslash{}input\{preamble\}, so no preamble text lives inside the chapter itself. The card takes that file's lines 8-18 (the verbatim package and the three \textbackslash{}comment-based environments reference/slogan/history at lines 9-12, section 13.2, together with the Xy-pic setup) and its lines 154 to the runtime end of the file (the native theorem declarations and mathematical macros). Not one character was changed, no package was added, no macro was defined, and the excerpt serves only as preamble, never as a span coordinate.\par\medskip
\noindent\textit{Editorial note.} The same-chapter pointers are located at their \textbackslash{}label lines inside the pinned chapter file; the cross-chapter pointers are located in the project's own tag list staged in the snapshot (sources/stacks-chow/tags.txt), at the verbatim lines carrying their canonical tags.\par\medskip
\noindent\textit{Editorial note.} The counted proof refers to two same-chapter lemmas by label; both are recorded as navigation stubs pointed at their \textbackslash{}label lines, and the arguments consuming them stay inside the counted span.\par\medskip
\noindent\textit{Editorial note.} The counted spans contain no \textbackslash{}cite, so the card carries no bibliography.\par\medskip
\noindent\textit{Allowed standard prerequisites (editorial).}\begin{itemize}
\item Radical of an ideal and its behaviour under inverse image. Used at the start of the counted proof (line 6057). Shows the conditions on the radicals of extended ideals and on radical ideals are equivalent.
\item Membership in the image via nonvanishing after tensoring with the residue field. Invoked in the counted proof at lines 6064-6067 and again at the end. Connects the contraction condition for primes with the topological surjectivity statement.
\end{itemize}\medskip
\ifdefined\hypertarget\hypertarget{sourcepointer-1}{}\fi
\noindent\textit{Source reference.} Original lemma 00E0: sources/stacks-algebra/pinned-original-source.tex, source lines 3204--3204. Navigation pointer for \textbackslash{}ref\{lemma-Zariski-topology\} used inside the counted span (line 6062); the argument that consumes it is retained in the span.\par\medskip
\ifdefined\hypertarget\hypertarget{sourcepointer-2}{}\fi
\noindent\textit{Source reference.} Original lemma 00E7: sources/stacks-algebra/pinned-original-source.tex, source lines 3717--3717. Navigation pointer for \textbackslash{}ref\{lemma-in-image\} used inside the counted span (line 6065); the argument that consumes it is retained in the span.\par\medskip
\begin{lemma}
\label{lemma-surjective-spec-radical-ideal}
Let $\varphi : R \to S$ be a ring map.
The following are equivalent:
\begin{enumerate}
\item The map $\Spec(S) \to \Spec(R)$ is surjective.
\item For any ideal $I \subset R$
the inverse image of $\sqrt{IS}$ in $R$ is equal to $\sqrt{I}$.
\item For any radical ideal $I \subset R$ the inverse image
of $IS$ in $R$ is equal to $I$.
\item For every prime $\mathfrak p$ of $R$ the inverse
image of $\mathfrak p S$ in $R$ is $\mathfrak p$.
\end{enumerate}
In this case the same is true after any base change: Given a ring map
$R \to R'$ the ring map $R' \to R' \otimes_R S$ has the equivalent
properties (1), (2), (3) as well.
\end{lemma}

\begin{proof}
If $J \subset S$ is an ideal, then
$\sqrt{\varphi^{-1}(J)} = \varphi^{-1}(\sqrt{J})$. This shows that (2)
and (3) are equivalent.
The implication (3) $\Rightarrow$ (4) is immediate.
If $I \subset R$ is a radical ideal, then
Lemma \ref{lemma-Zariski-topology}
guarantees that $I = \bigcap_{I \subset \mathfrak p} \mathfrak p$.
Hence (4) $\Rightarrow$ (2). By
Lemma \ref{lemma-in-image}
we have $\mathfrak p = \varphi^{-1}(\mathfrak p S)$ if and only if
$\mathfrak p$ is in the image. Hence (1) $\Leftrightarrow$ (4).
Thus (1), (2), (3), and (4) are equivalent.

\medskip\noindent
Assume (1) holds. Let $R \to R'$ be a ring map. Let
$\mathfrak p' \subset R'$ be a prime ideal lying over the prime
$\mathfrak p$ of $R$. To see that $\mathfrak p'$ is in the image
of $\Spec(R' \otimes_R S) \to \Spec(R')$ we have to show
that $(R' \otimes_R S) \otimes_{R'} \kappa(\mathfrak p')$ is not zero, see
Lemma \ref{lemma-in-image}.
But we have
$$
(R' \otimes_R S) \otimes_{R'} \kappa(\mathfrak p') =
S \otimes_R \kappa(\mathfrak p)
\otimes_{\kappa(\mathfrak p)} \kappa(\mathfrak p')
$$
which is not zero as $S \otimes_R \kappa(\mathfrak p)$ is not zero
by assumption and $\kappa(\mathfrak p) \to \kappa(\mathfrak p')$ is
an extension of fields.
\end{proof}

\end{document}
